\documentclass[11pt]{article}
\usepackage[margin=1in]{geometry}
\usepackage{amsmath,amssymb,amsthm}
\usepackage{graphicx}
\usepackage{hyperref}
\graphicspath{{runs/fa_banach_001/solutions/partial/2211.01544_Bx_weakly_null_c0_border/}}

\newtheorem{theorem}{Theorem}
\newtheorem{corollary}{Corollary}
\newtheorem{remark}{Remark}

\newcommand{\N}{\mathbb N}
\newcommand{\Rideal}{\mathcal R}
\newcommand{\Bideal}{\mathcal B}

\title{A norm-null/\(c_0\) border for Question 5.15 of Martínez--Meza-Alcántara--Uzcátegui}
\author{fa\_banach\_001, agent\_lane\_07}
\date{2026-06-20}

\begin{document}
\maketitle

\section*{Source question}

The source is Jorge Martínez, David Meza-Alcántara, and Carlos Uzcátegui, \emph{Pathology of submeasures and \(F_\sigma\) ideals}, arXiv:2211.01544v2.  After proving that, for sequences in \(c_0\),
\[
        \Rideal\le_K \Bideal({\bf x}) \quad\Longleftrightarrow\quad (x_n)_n
        \text{ is weakly null},
\]
the authors ask:

\begin{quote}
Which Banach spaces satisfy that \(\Rideal\le_K\Bideal({\bf x})\) for every weakly null sequence \({\bf x}=(x_n)_n\)?
\end{quote}

\begin{center}
\includegraphics[width=0.96\linewidth]{figures/open_problem_crop.png}
\end{center}

\section*{Idea of the proof}

There are two positive mechanisms.  First, if \(\|x_n\|\to0\), then the submeasure associated to \(\Bideal({\bf x})\) has property A in the sense of Section 4 of the source paper.  Proposition 4.8 of the source then gives
\(\Rideal\le_K \Bideal({\bf x})\).  Thus all Schur spaces satisfy the question.

Second, the source theorem already solves the problem for \(c_0\).  This is stable under passing to subspaces of \(c_0\), because the ideal \(\Bideal({\bf x})\) only records uniform boundedness of finite partial sums, and this is preserved by an isomorphic embedding.

The negative part uses the opposite mechanism.  If \(\Rideal\le_K\Bideal({\bf x})\), then \(\Bideal({\bf x})\) must be tall.  For a seminormalized weakly null sequence with no \(c_0\)-subsequence, no infinite subseries can be perfectly bounded: otherwise it would be weakly unconditionally Cauchy, and the Bessaga-Pelczynski selection principle would produce a \(c_0\)-subsequence.

\section*{Partial answer}

\begin{theorem}
\label{thm:border}
Let \(X\) be a Banach space.
\begin{enumerate}
\item For every sequence \({\bf x}=(x_n)_n\) in any Banach space, if \(\|x_n\|\to0\), then
\[
        \Rideal\le_K \Bideal({\bf x}).
\]
Consequently every Schur space satisfies Question 5.15.
\item If \(X\) is isomorphic to a subspace of \(c_0\), then for every weakly null sequence \({\bf x}=(x_n)_n\) in \(X\) one has
\[
        \Rideal\le_K \Bideal({\bf x}).
\]
\item If \(X\) contains a seminormalized weakly null sequence with no subsequence equivalent to the unit vector basis of \(c_0\), then \(X\) does not satisfy the property in Question 5.15.
\end{enumerate}
\end{theorem}

\begin{proof}
For (1), let \(\varphi_{\bf x}\) be the lower semicontinuous submeasure associated to \({\bf x}\) in Lemma 5.2 of the source paper:
\[
        \varphi_{\bf x}(A)=
        \sup\left\{\left\|\sum_{n\in F}x_n\right\|:
        \varnothing\ne F\in[A]^{<\omega}\right\}.
\]
Lemma 5.2 gives \(\Bideal({\bf x})=\mathrm{FIN}(\varphi_{\bf x})\).  If \(\|x_n\|\to0\), then for every \(\varepsilon>0\) the set
\[
        M_\varepsilon=\{n:\varphi_{\bf x}(\{n\})>\varepsilon\}
        =\{n:\|x_n\|>\varepsilon\}
\]
is finite, hence belongs to \(\mathrm{FIN}(\varphi_{\bf x})\).  Thus \(\varphi_{\bf x}\) has property A.  Proposition 4.8 of the source paper gives
\[
        \Rideal\le_K\mathrm{FIN}(\varphi_{\bf x})=\Bideal({\bf x}).
\]
If \(X\) has the Schur property, every weakly null sequence in \(X\) is norm-null, so the preceding argument applies to every weakly null sequence.

For (2), let \(J:X\to c_0\) be an isomorphic embedding and let \((x_n)_n\) be weakly null in \(X\).  Then \((Jx_n)_n\) is weakly null in \(c_0\), since every coordinate functional on \(c_0\), composed with \(J\), is a bounded functional on \(X\).  By Theorem 5.14 of the source paper,
\[
        \Rideal\le_K\Bideal((Jx_n)_n).
\]
It remains only to identify the ideal.  Since \(J\) is an isomorphic embedding, there are constants \(a,b>0\) such that
\[
        a\|u\|_X\le \|Ju\|_{c_0}\le b\|u\|_X\qquad (u\in X).
\]
Therefore, for every \(A\subseteq\N\),
\[
        \sup_{F\in[A]^{<\omega}}\left\|\sum_{n\in F}x_n\right\|_X<\infty
        \quad\Longleftrightarrow\quad
        \sup_{F\in[A]^{<\omega}}\left\|\sum_{n\in F}Jx_n\right\|_{c_0}<\infty.
\]
Thus \(\Bideal({\bf x})=\Bideal((Jx_n)_n)\), and the desired Katetov reduction follows.

For (3), let \((x_n)_n\) be seminormalized, weakly null, and have no \(c_0\)-subsequence.  We claim that \(\Bideal({\bf x})\) is not tall.  Indeed, if some infinite \(A\subseteq\N\) belonged to \(\Bideal({\bf x})\), then the series \(\sum_{n\in A}x_n\) would be perfectly bounded.  By Proposition 5.12 of the source paper, equivalently,
\[
        \sum_{n\in A}|x^*(x_n)|<\infty
        \qquad\text{for every }x^*\in X^*,
\]
so the subseries is weakly unconditionally Cauchy.  Since \((x_n)_{n\in A}\) is seminormalized, the classical Bessaga-Pelczynski \(c_0\) principle yields a subsequence equivalent to the unit vector basis of \(c_0\), contradicting the choice of \((x_n)\).  Hence \(\Bideal({\bf x})\) contains no infinite subset of \(\N\), so it is not tall.

The source paper recalls that \(\Rideal\le_K\mathcal I\) implies \(\mathcal I\) is tall.  Consequently \(\Rideal\not\le_K\Bideal({\bf x})\), and \(X\) fails Question 5.15.
\end{proof}

\begin{corollary}
Every infinite-dimensional reflexive Banach space fails the property in Question 5.15.
\end{corollary}

\begin{proof}
Every infinite-dimensional reflexive Banach space contains a seminormalized weakly null basic sequence.  Since a reflexive space contains no copy of \(c_0\), that sequence has no \(c_0\)-subsequence.  Apply Theorem~\ref{thm:border}(2).
\end{proof}

\section*{Verification and limitations}

This is a partial answer, not a full characterization.  It gives two broad positive classes, Schur spaces and subspaces of \(c_0\), and a broad obstruction, all spaces containing a seminormalized weakly null sequence without a \(c_0\)-subsequence.  The remaining border concerns non-Schur spaces where every seminormalized weakly null sequence has a \(c_0\)-subsequence, but which are not known or not assumed to embed into \(c_0\).  At the sequence level, the issue is that tallness of \(\Bideal({\bf x})\) is weaker than \(\Rideal\le_K\Bideal({\bf x})\); the source paper's Question 4.17 asks whether every nonpathological tall \(F_\sigma\) ideal is above \(\Rideal\), so this frontier is tied to the paper's broader Katetov problem.

Bounded novelty checks used the run indexes for `2211.01544`, `B(x)`, weakly null, random ideal, Schur property, and \(c_0\), plus web searches for the exact wording of Question 5.15 and related phrases.  No later explicit answer was found.  The nearby run packet on arXiv:2508.21235 concerns critical ideals of ideal Schauder bases and nonpathological Banach-representable ideals; it is related background but does not record this norm-null/\(c_0\)-border statement for Question 5.15.

\begin{thebibliography}{9}
\bibitem{MMU}
J. Martínez, D. Meza-Alcántara, and C. Uzcátegui, \emph{Pathology of submeasures and \(F_\sigma\) ideals}, arXiv:2211.01544.

\bibitem{AK}
F. Albiac and N. J. Kalton, \emph{Topics in Banach Space Theory}, Springer, 2016.
\end{thebibliography}

\end{document}
